\chapter{Réalisation}
\label{chap:realisation}

\section{Introduction}
\label{sec:re-intro}

Sur la base des choix arrêtés au chapitre précédent, ce chapitre présente la mise en \oe{}uvre concrète de l'architecture proposée. Nous décrivons successivement~: l'environnement de développement, l'architecture globale et la topologie réseau, la PKI Vault, le broker MQTT en TLS/mTLS, les scénarios d'attaque retenus avec leurs contre-mesures, le pipeline d'apprentissage automatique (présenté \emph{après} la couche de défense statique pour respecter une logique de défense en profondeur), puis les captures d'écran des interfaces clés (Vault UI, Wazuh, GNS3). L'objectif est de montrer comment chaque brique conçue a été effectivement instanciée dans le laboratoire expérimental.

\section{Environnement de développement}
\label{sec:re-env}

Le laboratoire est entièrement reproductible sur une station Windows~10/11 (16~Go de RAM minimum). Les principaux composants logiciels sont résumés dans le tableau~\ref{tab:env}.

\begin{table}[!htbp]
\centering
\caption{Composants de l'environnement de développement}
\label{tab:env}
\small
\begin{tabularx}{\linewidth}{|l|l|X|}
\hline
\textbf{Composant} & \textbf{Version} & \textbf{Rôle} \\
\hline
GNS3 & 2.2.x & Émulation de la topologie réseau \\ \hline
Docker Engine & 24.x & Conteneurs services et clients IoT \\ \hline
HashiCorp Vault & 1.15 & PKI à deux niveaux et secrets engine \\ \hline
Eclipse Mosquitto & 2.0 & Broker MQTT avec TLS~1.3 + mTLS \\ \hline
Suricata & 7.0 & IDS réseau, parser MQTT \\ \hline
Wazuh & 4.7 & SIEM, agrégation et corrélation \\ \hline
pfSense & 2.7 & Pare-feu inter-VLAN \\ \hline
MikroTik RouterOS & 7.x & Routage / NAT \\ \hline
Python & 3.11 & Simulateur de flotte, ML \\ \hline
scikit-learn & 1.4 & Algorithmes ML \\ \hline
\end{tabularx}
\end{table}

Ces composants ne fonctionnent pas de manière indépendante~: ils s'articulent en une \textbf{architecture de défense en profondeur} où chaque couche renforce les autres. La section suivante décrit cette organisation globale et le flux de données de bout en bout, avant d'entrer dans le détail de chaque brique.

\section{Architecture globale et topologie réseau}
\label{sec:re-archi}

\subsection{Topologie réseau}

La figure~\ref{fig:topo-p1} présente la topologie réseau du laboratoire~: la chaîne Internet~\textrightarrow{} pfSense~\textrightarrow{} MikroTik et la segmentation en quatre VLANs, avec un miroir passif du trafic DMZ vers Suricata. L'architecture générale du système a déjà été introduite au chapitre~\ref{chap:introduction} (figure~\ref{fig:archi-reelle-ch1})~; la présente section en détaille la topologie réseau effective.

\begin{figure}[!htbp]
\centering
\includegraphics[width=0.90\linewidth]{diag-topo-part1.png}
\caption{Topologie réseau~: infrastructure pfSense/MikroTik et segmentation en 4~VLANs (IoT, DMZ, PKI, SIEM)}
\label{fig:topo-p1}
\end{figure}
\FloatBarrier

La segmentation repose sur quatre zones isolées~:

\begin{itemize}[leftmargin=*]
  \item \textbf{VLAN~10 (Zone IoT)}~: héberge le simulateur de flotte et les 50~dispositifs émulés. Les seuls flux sortants autorisés sont IoT\textrightarrow{}Broker/8883 (MQTT/TLS) et IoT\textrightarrow{}Vault/8200 (émission de certificats)~;
  \item \textbf{VLAN~20 (Zone DMZ)}~: contient Mosquitto (broker MQTT) et Suricata (IDS). Un port miroir MikroTik copie passivement le trafic de cette zone vers Suricata~;
  \item \textbf{VLAN~30 (Zone PKI)}~: isole HashiCorp Vault et la CA intermédiaire. Seul le flux PKI\textrightarrow{}SIEM/1514 est autorisé en sortie~;
  \item \textbf{VLAN~40 (Zone SIEM)}~: regroupe Wazuh, OpenSearch et le moteur ML. Les journaux de toutes les zones convergent ici via Syslog/1514.
\end{itemize}

La politique de filtrage inter-VLAN suit le principe \emph{deny-by-default}~: seuls les flux explicitement listés ci-dessus sont autorisés par pfSense, tout autre trafic inter-zones est bloqué.

Le cycle de fonctionnement de bout en bout suit cinq phases successives~: (1)~provisionnement de l'identité cryptographique via Vault~; (2)~établissement de la connexion TLS~1.3 mTLS avec le broker~; (3)~publication des données IoT sur les topics autorisés~; (4)~détection et supervision par Suricata et Wazuh~; (5)~rotation automatique des certificats à l'approche de leur expiration. Les deux phases les plus critiques~--- l'établissement du mTLS et la détection~--- sont détaillées dans les figures suivantes.

\begin{figure}[!htbp]
\centering
\includegraphics[width=0.95\linewidth,height=0.40\textheight,keepaspectratio]{diag-flux-phase2.png}
\caption{Flux E2E --- Phase~2~: établissement de la connexion TLS~1.3 mTLS entre un dispositif et le broker}
\label{fig:flux-p2}
\end{figure}
\FloatBarrier

\begin{figure}[!htbp]
\centering
\includegraphics[width=0.92\linewidth,height=0.38\textheight,keepaspectratio]{diag-flux-phase4.png}
\caption{Flux E2E --- Phase~4~: détection et supervision (Suricata IDS~\textrightarrow{} Wazuh SIEM~\textrightarrow{} tableau de bord)}
\label{fig:flux-p4}
\end{figure}
\FloatBarrier

\section{Infrastructure à clés publiques (Vault)}
\label{sec:re-pki}

\subsection{Hiérarchie}

La PKI suit le modèle à deux niveaux préconisé par le NIST~\cite{nist_sp800213}~: une CA racine \emph{offline} (générée hors ligne, conservée chiffrée) signe une unique CA intermédiaire en ligne, gérée par Vault~\cite{vault_pki}. La figure~\ref{fig:pki-hier} illustre cette hiérarchie.

\begin{figure}[!htbp]
\centering
\includegraphics[width=0.85\linewidth]{diag-pki-hierarchy.png}
\caption{Hiérarchie de la PKI à deux niveaux (CA racine \emph{offline} + CA intermédiaire Vault)}
\label{fig:pki-hier}
\end{figure}

Le diagramme distingue les trois niveaux de confiance de la hiérarchie~: (i)~la \textbf{CA racine} (RSA-4096, self-signed, TTL~10~ans) est conservée hors ligne et ne signe que les CA intermédiaires~; sa clé privée n'est jamais exposée dans la VM GNS3~; (ii)~la \textbf{CA intermédiaire} (RSA-2048, TTL~5~ans) est embarquée dans le \emph{secrets engine PKI} de Vault et signe les certificats opérationnels à la demande via API REST~; (iii)~les \textbf{certificats finaux} sont émis selon deux rôles distincts~: \texttt{role-server} (TTL~90~jours, usage \texttt{serverAuth}, pour Mosquitto et Vault), et \texttt{role-device} (TTL~30~jours, usage \texttt{clientAuth}, CN contrôlé au pattern \texttt{dev-XXX.iot.lab}). La séparation des rôles empêche qu'un certificat de service soit utilisé pour s'authentifier comme dispositif, et inversement.

\subsection{Rôles Vault et émission}

Deux rôles PKI ont été définis~: \texttt{role-server} (TTL 90~jours, usage \emph{server\_auth}) et \texttt{role-device} (TTL 30~jours, usage \emph{client\_auth}, \emph{Common Name} contraint au pattern \texttt{dev-XXX.iot.lab}). L'émission d'un certificat client se réduit à un appel HTTP authentifié par jeton AppRole.

\begin{lstlisting}[language=bash, caption={Émission d'un certificat client via Vault}]
vault write -format=json pki_int/issue/role-device \
    common_name="dev-042.iot.lab" \
    ttl="720h"
\end{lstlisting}

Le choix des TTL résulte d'une analyse du compromis \textbf{sécurité / coût opérationnel}~: un TTL court réduit la fenêtre d'exposition en cas de vol de clé privée, mais impose une rotation fréquente, source de charge et d'erreur. Pour les services (broker, Vault), un TTL de 90~jours équilibre la stabilité des configurations et la rotation régulière~; pour les dispositifs, un TTL de 30~jours est préférable car ces objets ont un cycle de vie moins maîtrisé et la rotation est automatisée par Vault, rendant son surcoût négligeable. Ce choix est cohérent avec les recommandations du NIST~\cite{nist_sp800213} qui préconise des TTL inférieurs à 90~jours pour les certificats d'entité terminale dans les environnements IoT.

La PKI étant opérationnelle, chaque composant~-- dispositif comme service~-- dispose d'une identité cryptographique vérifiable par tous les pairs. C'est sur cette fondation que repose la couche transport~: le broker MQTT peut désormais exiger la présentation d'un certificat signé par la CA intermédiaire avant d'accepter toute connexion, comme décrit dans la section suivante.

\section{Broker MQTT en TLS 1.3 + mTLS}
\label{sec:re-mqtt}

Mosquitto est configuré~\cite{mosquitto_tls} pour exiger un certificat client, associer l'identité certifiée au nom d'utilisateur et appliquer un fichier d'\acs{ACL} liant chaque \emph{Common Name} à ses topics autorisés. La figure~\ref{fig:tls-handshake} retrace le \emph{handshake} TLS~1.3 mutuel observé entre un dispositif et le broker.

\begin{figure}[!htbp]
\centering
\includegraphics[width=0.9\linewidth]{diag-tls-handshake.png}
\caption{\emph{Handshake} TLS~1.3 avec authentification mutuelle (mTLS)}
\label{fig:tls-handshake}
\end{figure}

Le diagramme de séquence se lit selon quatre phases distinctes~:
\begin{enumerate}[leftmargin=*]
  \item \textbf{ClientHello} (dispositif~\textrightarrow~broker)~: le dispositif annonce les suites cryptographiques qu'il supporte et propose les paramètres de partage de clé (ECDHE)~;
  \item \textbf{ServerHello + CertificateVerify} (broker~\textrightarrow~dispositif)~: le broker sélectionne la suite, transmet son certificat X.509 signé par la CA Vault et prouve la possession de sa clé privée~;
  \item \textbf{Certificate (client)} (dispositif~\textrightarrow~broker)~: le dispositif présente son propre certificat (rôle \texttt{role-device}, CN \texttt{dev-XXX.iot.lab})~; le broker vérifie sa validité, son appartenance à la chaîne PKI et l'absence de révocation dans la CRL Vault~;
  \item \textbf{Finished + Application Data}~: les deux parties dérivent les clés de session symétriques (HKDF sur secret ECDHE) et la session chiffrée est établie en \textbf{1-RTT}.
\end{enumerate}
L'authentification mutuelle garantit qu'aucun dispositif sans certificat valide ne peut accéder au broker, et qu'aucun broker frauduleux ne peut se faire passer pour Mosquitto.

L'architecture de confiance est ainsi opérationnelle~: les identités sont gérées par Vault, le canal de communication est chiffré et les accès applicatifs sont restreints par ACL par certificat. Pour valider la robustesse de cet ensemble, six scénarios d'attaque ont été conçus en s'appuyant directement sur les vecteurs identifiés dans l'état de l'art~: interception, usurpation d'identité, déni de service et exfiltration applicative.

\section{Scénarios d'attaque et contre-mesures}
\label{sec:re-attaques}

Six scénarios ont été conçus pour couvrir les principales catégories d'attaques applicables à un broker MQTT. Chaque scénario est défini par un \emph{vecteur}, une \emph{méthode de détection} et une \emph{contre-mesure}. Les scénarios sont rejoués depuis la flotte simulée et corroborés par les alertes Suricata/Wazuh.

\begin{attackbox}[Scénario 1 -- Man-in-the-Middle (TLS désactivé)]
\textbf{Vecteur~:} un dispositif tente de se connecter sur le port 1883 (MQTT en clair) au lieu de 8883 (TLS). Un attaquant en position MITM peut alors capturer puis modifier les messages.\\
\textbf{Détection~:} règle Suricata sur tout flux TCP/1883.\\
\textbf{Contre-mesure~:} pare-feu pfSense bloquant 1883~; broker écoutant uniquement sur 8883 avec \texttt{require\_certificate}.
\end{attackbox}

\begin{attackbox}[Scénario 2 -- Sniffing MQTT en clair]
\textbf{Vecteur~:} interception passive du trafic MQTT non chiffré sur le segment IoT.\\
\textbf{Détection~:} alerte Suricata si un \texttt{PUBLISH} apparaît hors de la session TLS.\\
\textbf{Contre-mesure~:} TLS~1.3 obligatoire. Les suites cryptographiques sont restreintes conformément à la RFC~9325~\cite{rfc9325}.
\end{attackbox}

\begin{attackbox}[Scénario 3 -- Brute-force d'authentification]
\textbf{Vecteur~:} 500~tentatives de \texttt{CONNECT} en moins de 60~s avec des certificats forgés.\\
\textbf{Détection~:} règle Suricata MQTT (seuil de \texttt{CONNECT} par IP source).\\
\textbf{Contre-mesure~:} certificats clients signés par CA Vault~; \emph{rate-limit} pfSense sur le port 8883.
\end{attackbox}

\begin{attackbox}[Scénario 4 -- Exfiltration via topic non autorisé]
\textbf{Vecteur~:} un dispositif compromis tente de \texttt{SUBSCRIBE} à \texttt{tt/admin/\#}.\\
\textbf{Détection~:} log Mosquitto \texttt{ACL denied}~; règle Wazuh corrélant N occurrences/15~min.\\
\textbf{Contre-mesure~:} ACL stricte par \emph{Common Name}, granularité au topic.
\end{attackbox}

\begin{attackbox}[Scénario 5 -- Déni de service par flood CONNECT]
\textbf{Vecteur~:} client malveillant ouvrant plusieurs milliers de connexions TLS partielles.\\
\textbf{Détection~:} alertes Suricata sur taux d'ouvertures TLS anormal.\\
\textbf{Contre-mesure~:} \texttt{max\_connections} Mosquitto, \texttt{state-table-limit} pfSense, blacklist dynamique Wazuh.
\end{attackbox}

\begin{attackbox}[Scénario 6 -- Compromission de certificat client]
\textbf{Vecteur~:} certificat \texttt{dev-017} volé puis réutilisé depuis une IP non habituelle.\\
\textbf{Détection~:} module ML (Isolation Forest) signalant l'écart comportemental~; corrélation Wazuh IP$\leftrightarrow$identité.\\
\textbf{Contre-mesure~:} révocation Vault (\texttt{vault write pki\_int/revoke serial=...}) et \emph{rotation} immédiate de la CRL.
\end{attackbox}

La chaîne de détection fonctionne en trois étapes. Suricata reçoit d'abord une copie passive du trafic VLAN~20 via le port miroir MikroTik (AF\_PACKET, zéro copie, aucun impact sur la latence MQTT), décode la pile réseau et active son parseur MQTT sur le port~8883 pour extraire les métadonnées applicatives (type de paquet, topic, QoS, taille de payload, \texttt{client\_id}). Les règles SID personnalisées sont ensuite appliquées, comme détaillé ci-dessous.

\begin{figure}[!htbp]
\centering
\includegraphics[width=0.73\linewidth,height=0.72\textheight,keepaspectratio]{diag-suricata-part2.png}
\caption{Chaîne de détection Suricata~\textrightarrow{}Wazuh~: application des règles SID personnalisées et génération d'alertes}
\label{fig:detection-p2}
\end{figure}
\FloatBarrier

La figure~\ref{fig:detection-p2} détaille les quatre règles SID personnalisées~: \textbf{SID~9000001} (flood~: >~50 CONNECT en 10~s depuis la même IP, 23~déclenchements)~; \textbf{SID~9000002} (topic suspect \texttt{cmd/}, 4~déclenchements)~; \textbf{SID~9000003} (payload~>~4~096~octets, 15~déclenchements)~; \textbf{SID~9000004} (scan broker~: >~20 CONNECT en 5~s sans SUBSCRIBE, 9~déclenchements). Total~: 51~alertes sur 76~000~messages, 0~faux positif sur le trafic TLS normal.

\begin{figure}[!htbp]
\centering
\includegraphics[width=0.73\linewidth,height=0.72\textheight,keepaspectratio]{diag-suricata-part3.png}
\caption{Pipeline d'ingestion des alertes~: EVE JSON~\textrightarrow{} Filebeat~\textrightarrow{} Wazuh~\textrightarrow{} OpenSearch~\textrightarrow{} Dashboard}
\label{fig:detection-p3}
\end{figure}
\FloatBarrier

La figure~\ref{fig:detection-p3} décrit le pipeline d'ingestion~: Suricata écrit chaque alerte dans \texttt{eve.json} (champs~: \texttt{timestamp}, \texttt{src\_ip}, \texttt{alert.signature}, \texttt{mqtt.topic}), Filebeat le transmet à Wazuh via Syslog/1514. Wazuh applique ses décodeurs XML puis ses règles de corrélation (ex.~: >~5~alertes SID~9000001 en~<~60~s → escalade niveau~12 «~MQTT DoS flood confirmé~»). OpenSearch indexe les événements et le tableau de bord les restitue à l'analyste SOC avec filtrage par type d'attaque, VLAN source et période. La latence bout-en-bout est inférieure à 2~secondes.

\section{Pipeline d'apprentissage automatique}
\label{sec:re-ml}

Les scénarios d'attaque ont mis en évidence une limite inhérente aux approches par règles~: le scénario~S6 (certificat valide réutilisé depuis un contexte inhabituel) n'est détecté que partiellement par Suricata et Wazuh. Le pipeline ML intervient précisément pour combler cette lacune~: il complète la détection par signature en repérant les comportements anormaux que les règles statiques ne peuvent pas anticiper. Il est constitué de trois phases~: préparation des données, entraînement des modèles et intégration dans le SIEM.

\paragraph{Phase~1 — Chargement, exploration et feature engineering.}
Le dataset CSV (76~000~entrées, 24~features brutes) est chargé puis soumis à une analyse exploratoire~: distribution des classes, matrice de corrélation et boxplots par type d'attaque. Les variables catégorielles (\texttt{protocol}, \texttt{tls\_version}, \texttt{cert\_status}, etc.) sont encodées via \texttt{LabelEncoder} pour obtenir 29~features exploitables. Deux variables cibles sont construites~: \texttt{label\_secure} (binaire~: 0~=~attaque, 1~=~normal) et \texttt{attack\_type} (multi-classe~: 7~types d'attaques).

\paragraph{Phase~2 — Entraînement des trois modèles.}
Les trois modèles sont entraînés en parallèle sur les mêmes 29~features (figure~\ref{fig:ml-p2}).

\begin{figure}[!htbp]
\centering
\includegraphics[width=0.80\linewidth]{diag-ml-part2.png}
\caption{Pipeline ML~: entraînement parallèle des trois modèles (IsolationForest, RF binaire, RF multi-classes)}
\label{fig:ml-p2}
\end{figure}
\FloatBarrier

\begin{itemize}[leftmargin=*]
  \item \textbf{Modèle~A — IsolationForest} (non supervisé)~: entraîné exclusivement sur le trafic normal (\texttt{label\_secure~=~1}), il calcule un score d'anomalie pour chaque message, avec un seuil fixé à $-0{,}05$. Résultats~: F1~=~0,932 et ROC-AUC~=~0,959~;
  \item \textbf{Modèle~B — RandomForest binaire} (supervisé)~: entraîné avec un split stratifié 80/20, \texttt{class\_weight=balanced} et \texttt{n\_estimators=200}. Résultats~: F1~=~0,994 et ROC-AUC~=~0,999~;
  \item \textbf{Modèle~C — RandomForest multi-classes}~: classification directe des 7~types d'attaques, matrice de confusion sans confusion Normal$\leftrightarrow$DoS. Résultats~: F1~macro~=~0,993 et ROC-AUC~=~0,998.
\end{itemize}

\paragraph{Phase~3 — Comparaison, export et intégration SIEM.}
Les trois modèles sont comparés sur F1, ROC-AUC, précision et rappel. La stratégie retenue est bicouche~: IsolationForest pour la détection sans étiquettes (données de production non labellisées) et RandomForest binaire pour la classification en temps réel du trafic labelé. Les modèles sont sauvegardés au format \texttt{.pkl} (joblib) dans \texttt{ml/models/}. Les scores d'anomalie sont transmis à Wazuh via l'API REST (port~55000) et corrélés avec les alertes Suricata pour une détection multicouche~: un message simultanément signalé par l'IDS et classé anormal par le modèle ML déclenche une alarme de niveau critique.

\subsection{Données et features}

Le jeu de données utilisé contient 76\,000 enregistrements répartis en huit classes (1 normale + 7 attaques)~\cite{hasan2020securing}. Les figures~\ref{fig:eda-classes}--\ref{fig:eda-corr} présentent l'analyse exploratoire.

\begin{figure}[!htbp]
\centering
\includegraphics[width=0.85\linewidth,height=0.38\textheight,keepaspectratio]{eda_classes_distribution.png}
\caption{Analyse exploratoire --- distribution des classes}
\label{fig:eda-classes}
\end{figure}
\FloatBarrier

Ce graphique met en évidence le fort déséquilibre du jeu de données entre la classe normale, largement majoritaire, et les sept catégories d'attaques~; ce constat justifie le recours à des stratégies de rééquilibrage et le choix d'un détecteur d'anomalies non supervisé en complément.

\begin{figure}[!htbp]
\centering
\includegraphics[width=0.95\linewidth,height=0.45\textheight,keepaspectratio]{eda_correlation_heatmap.png}
\caption{Analyse exploratoire --- matrice de corrélation des features numériques}
\label{fig:eda-corr}
\end{figure}
\FloatBarrier

La matrice de corrélation révèle les variables fortement liées entre elles, ce qui permet d'écarter les features redondantes et de retenir un sous-ensemble informatif pour l'entraînement des modèles.
\FloatBarrier

\paragraph{Synthèse de présentation.}
Les interfaces d'administration (Vault, Wazuh, GNS3) ont été mobilisées tout au long du déploiement et de la validation, mais leurs captures détaillées ne sont pas conservées dans le corps principal afin de privilégier les choix d'architecture, les mécanismes de sécurité et les résultats mesurés. Les éléments visuels les plus probants sont repris au chapitre~\ref{chap:tests-conclusion}, sous forme de preuves ciblées directement liées aux scénarios évalués.

\section{Conclusion}
\label{sec:re-conclusion}

Ce chapitre a décrit la mise en \oe{}uvre concrète de l'architecture conçue~: un environnement reproductible sous GNS3 et Docker, une PKI Vault à deux niveaux assurant l'émission et la révocation automatisées, un broker Mosquitto imposant mTLS à toutes les connexions, une flotte nominale de 50~dispositifs rejouant un jeu de données de 76\,000~messages, six scénarios d'attaque couvrant les principaux vecteurs IoT, et un pipeline ML complémentaire. Chacune de ces briques répond directement à l'un des objectifs O1--O7 définis au chapitre~\ref{chap:introduction}~: la réalisation valide ainsi la conception dans ses dimensions fonctionnelles et sécuritaires. Il reste désormais à \textbf{mesurer} l'efficacité, la performance et la portée expérimentale de cet ensemble~: c'est l'objet du chapitre suivant.
